\section{Lecture 7: Linear Regression and Geometry}

\subsection{Motivation}

Linear regression is one of the simplest yet most fundamental models in machine learning. Despite its apparent simplicity, it illustrates many core ideas:

\begin{itemize}
    \item function approximation,
    \item optimization,
    \item geometry of solutions,
    \item regularization and generalization.
\end{itemize}

\medskip

\noindent
\textbf{Guiding idea.}  
Linear regression provides a setting where learning can be understood exactly, both algebraically and geometrically.

\subsection{Linear Models and Projections}

We consider a linear model:
\[
f_w(x) = w^\top x, \quad w \in \mathbb{R}^d.
\]

\medskip

\noindent
Given data $(x_i, y_i)$, define the design matrix:
\[
X =
\begin{bmatrix}
x_1^\top \\
x_2^\top \\
\vdots \\
x_n^\top
\end{bmatrix}
\in \mathbb{R}^{n \times d}, \quad
y =
\begin{bmatrix}
y_1 \\
y_2 \\
\vdots \\
y_n
\end{bmatrix}.
\]

\medskip

\noindent
Predictions are given by:
\[
\hat{y} = Xw.
\]

\medskip

\noindent
\textbf{Geometric interpretation.}

\begin{itemize}
    \item The column space of $X$ is a subspace of $\mathbb{R}^n$
    \item Predictions $Xw$ lie in this subspace
\end{itemize}

\medskip

\noindent
\textbf{Projection viewpoint.}  
The optimal prediction is the projection of $y$ onto the column space of $X$.

\medskip

\noindent
\textbf{Orthogonality condition.}
\[
X^\top (Xw - y) = 0.
\]

\medskip

\noindent
\textbf{Insight.}  
Learning corresponds to projecting observed outputs onto the space of representable predictions.

\subsection{Least Squares as Optimization}

We minimize the squared loss:
\[
L(w) = \|Xw - y\|^2.
\]

\medskip

\noindent
\textbf{Normal equations.}
\[
X^\top X w = X^\top y.
\]

\medskip

\noindent
\textbf{Solution (full rank case).}
\[
w^* = (X^\top X)^{-1} X^\top y.
\]

\medskip

\noindent
\textbf{Interpretation.}
\begin{itemize}
    \item Closed-form solution exists
    \item Equivalent to minimizing empirical risk with squared loss
\end{itemize}

\medskip

\noindent
\textbf{Overparameterized case ($d > n$).}
\begin{itemize}
    \item Infinitely many solutions
    \item Minimum-norm solution:
    \[
    w^* = X^\top (XX^\top)^{-1} y
    \]
\end{itemize}

\medskip

\noindent
\textbf{Connection to optimization.}  
Gradient descent converges to this minimum-norm solution when initialized at $w=0$.

\medskip

\noindent
\textbf{Insight.}  
Even simple models reveal the role of implicit regularization.

\subsection{Conditioning and Numerical Stability}

The matrix $X^\top X$ plays a central role.

\medskip

\noindent
\textbf{Condition number.}
\[
\kappa = \frac{\lambda_{\max}}{\lambda_{\min}},
\]
where $\lambda_i$ are eigenvalues of $X^\top X$.

\medskip

\noindent
\textbf{Ill-conditioning.}
\begin{itemize}
    \item Small eigenvalues $\Rightarrow$ unstable inversion
    \item Solutions sensitive to noise
\end{itemize}

\medskip

\noindent
\textbf{Effects.}
\begin{itemize}
    \item Slow convergence of gradient descent
    \item Large variance in estimated parameters
\end{itemize}

\medskip

\noindent
\textbf{Geometric view.}
\begin{itemize}
    \item Level sets are elongated ellipsoids
    \item Optimization proceeds slowly along narrow directions
\end{itemize}

\medskip

\noindent
\textbf{Insight.}  
Numerical properties of the data matrix strongly affect learning behavior.

\subsection{Ridge Regression}

To address instability and overfitting, we introduce $\ell_2$ regularization:
\[
\min_w \|Xw - y\|^2 + \lambda \|w\|^2.
\]

\medskip

\noindent
\textbf{Solution.}
\[
w^* = (X^\top X + \lambda I)^{-1} X^\top y.
\]

\medskip

\noindent
\textbf{Interpretation.}
\begin{itemize}
    \item Adds stability to inversion
    \item Shrinks parameter values
\end{itemize}

\medskip

\noindent
\textbf{Geometric view.}
\begin{itemize}
    \item Solution lies in intersection of:
    \begin{itemize}
        \item loss level sets
        \item $\ell_2$ ball
    \end{itemize}
\end{itemize}

\medskip

\noindent
\textbf{Effect on eigenvalues.}
\[
\lambda_i \mapsto \lambda_i + \lambda,
\]
which improves conditioning.

\medskip

\noindent
\textbf{Bias--variance tradeoff.}
\begin{itemize}
    \item Increased bias
    \item Reduced variance
\end{itemize}

\medskip

\noindent
\textbf{Insight.}  
Ridge regression stabilizes learning and improves generalization.

\subsection{A Unifying Perspective}

Linear regression illustrates key principles of learning:

\begin{itemize}
    \item \textbf{Approximation:} linear hypothesis space
    \item \textbf{Optimization:} least squares objective
    \item \textbf{Geometry:} projection onto subspaces
    \item \textbf{Regularization:} explicit and implicit
\end{itemize}

\medskip

\noindent
\textbf{Core idea.}  
Learning is the projection of data onto a structured space, shaped by both the model and the optimization process.

\medskip

\noindent
\textbf{Key insight.}  
Even the simplest models reveal deep connections between geometry, optimization, and generalization.

\subsection{Outlook}

In this lecture, we developed a complete understanding of linear regression, including:

\begin{itemize}
    \item geometric interpretation via projections,
    \item closed-form solutions,
    \item numerical stability and conditioning,
    \item regularization through ridge regression.
\end{itemize}

\medskip

\noindent
These ideas extend naturally to classification problems.

\medskip

\noindent
In the next lecture, we study logistic regression, where probabilistic modeling and optimization come together in classification tasks.

\medskip

\noindent
\textbf{Guiding principle.}  
Simple models provide a foundation for understanding complex learning systems.